
\section{中药材烘干问题：两套独立方案的对比检验报告}



检验方法：<b>把两个方案的环境模型互换</b>，用两套完全不同的求解器交叉运行<br/>

全部结论均附可复现的脚本与原始输出



\section{第 1 章 检验对象与方法}


\subsection{两套方案}


\begin{center}
\small
\begin{tabular}{lll}
\toprule
 & 方案 X（「行」） & 方案 B（已审论文） \\
\midrule
控制方程 & $\rho c_p\partial_tT=r^{-1}\partial_r(rk\partial_rT)$；$\partial_tC=r^{-1}\partial_r(rD\partial_rC)$ & 同 \\
收缩处理 & 物质坐标 $\xi=r/R(t)$，由干物质守恒推出无对流项 & 材料坐标 $\sigma=r/R(t)$，由链式法则证明对流项精确抵消 \\
空间离散 & \textbf{节点中心有限体积}，$N=1600$ 个控制体 & \textbf{分级谱元弱形式}（Galerkin），ndof=55/81 \\
时间离散 & $\theta$ 法（CN）+ \textbf{Rannacher 启动 4 步} & 自适应隐式 BDF / Radau IIA \\
烘房环境 & \textbf{一阶惯性式拟合}（$T_{set}=50.212$，$\tau_T=1877.5$ s） & \textbf{分段线性插值}（14400 s 后取末端值 50.165） \\
终止判据 & 全场 $\max C<0.15$，在 60 s/10 s 网格上线性插值定位 & 连续事件函数 $C(0,t)-0.15=0$，Brent 求根 \\
非均匀收缩 & 未涉及 & 另建 M1 自洽模型（$-6.3\%$） \\
二维端面 & 用「轴向特征时间相差 156 倍」做量级论证 & 二维轴对称 + 同网格隔离实验（$-0.004\%$） \\
\bottomrule
\end{tabular}
\end{center}


\textbf{两套方案的物理模型完全相同}（同一组方程、同一组边界条件、同一组物性经验式、同一组题面参数），差别只在\textbf{离散方法}与\textbf{环境数据处理}。


\subsection{检验方法：环境互换实验}


要判断「两套结果的差异来自哪里」，最干净的做法是\textbf{只换一个因素}。我做了 2$\times$2 的交叉实验：


\begin{center}
\small
\begin{tabular}{lll}
\toprule
 & 方案 X 的环境（一阶惯性式） & 方案 B 的环境（分段线性） \\
\midrule
\textbf{方案 X 的求解器}（FV, N=1600, dt=1 s） & 实验 ① & 实验 ② \\
\textbf{方案 B 的求解器}（谱元, ndof=81, BDF） & 实验 ③ & 实验 ④ \\
\bottomrule
\end{tabular}
\end{center}


如果差异来自环境，那么 ① 与 ③ 应当一致、② 与 ④ 应当一致；如果差异来自离散方法，那么 ① 与 ② 应当一致、③ 与 ④ 应当一致。


\vspace{1mm}\hrule\vspace{1mm}


\section{第 2 章 核心结论}


\subsection{一句话结论}


\begin{quote}
\textbf{两套方案的差异 100\% 来自「如何把附件 1 的噪声数据变成边界条件」，与空间离散、时间离散、求解器实现完全无关。}  在同一环境下，\textbf{节点中心有限体积（1600 控制体）与分级谱元（81 自由度）在四位小数上完全一致}——两套数值方法互相验证通过。
\end{quote}


\subsection{交叉实验结果（问题 1，$t=1800$ s）}


\begin{center}
\small
\begin{tabular}{llllll}
\toprule
实验 & 求解器 & 环境 & $T_{中心}$ & $T_{表面}$ & $C_{表面}$ \\
\midrule
① & 方案 X FV & 方案 X & \textbf{34.0425} & \textbf{37.1989} & \textbf{1.5109} \\
② & 方案 X FV & 方案 B & \textbf{33.5753} & \textbf{36.7856} & \textbf{1.5102} \\
③ & 方案 B 谱元 & 方案 X & \textbf{34.0425} & \textbf{37.1989} & \textbf{1.5109} \\
④ & 方案 B 谱元 & 方案 B & \textbf{33.5753} & \textbf{36.7856} & \textbf{1.5102} \\
 & \textbf{方案 X 论文值} &  & 34.0425 & 37.1989 & 1.5109 \\
 & \textbf{方案 B 论文值} &  & 33.5753 & 36.7856 & 1.5102 \\
\bottomrule
\end{tabular}
\end{center}


\textbf{读法}：


\begin{itemize}
\item 同一行（同一环境）内，两种求解器给出一模一样的四位小数 ⟹ \textbf{离散方法无差异}；
\item 同一列（同一求解器）内，两种环境给出不同的结果 ⟹ \textbf{环境是唯一的差异源}；
\item ① = 方案 X 论文值、④ = 方案 B 论文值 ⟹ 两篇论文的数字都能被\textbf{另一方的求解器}精确复现。
\end{itemize}


\textbf{这条实验的价值}：它把「谁的数错了」这个问题，转换成了「两种环境处理哪种更合适」这个清晰、可讨论的建模选择问题。


\subsection{环境差异有多大}


\begin{center}
\small
\begin{tabular}{llll}
\toprule
$t$/s & 拟合式 $T_{air}$ & 分段线性 $T_{air}$ & 差 \\
\midrule
0 & 28.000 & 28.000 & 0.000 \\
300 & 31.280 & 30.834 & \textbf{+0.446} \\
600 & 34.076 & 33.202 & \textbf{+0.874} \\
1200 & 38.490 & 37.882 & +0.608 \\
1800 & 41.696 & 41.513 & +0.183 \\
3600 & 46.947 & 47.485 & \textbf{−0.538} \\
14400 & 50.212 & 50.165 & +0.047 \\
\bottomrule
\end{tabular}
\end{center}


拟合曲线在 \textbf{60–2300 s 内系统性高于实测}（最大 +0.96 ℃），在 2700–7200 s 内系统性低于实测。拟合残差的\textbf{滞后 1 阶自相关达 0.83}，说明这不是「去噪」，而是\textbf{模型形式失配}。


\textbf{方案 X 对此是知情且披露的}（原文 §2.2："拟合曲线在 $60\sim2300$ s 内系统性高于实测，最大偏高 $0.96\,^\circ$C……使表 1 各值系统性偏高，最大偏高 $0.52\,^\circ$C"）。这一点应当给予肯定。


\vspace{1mm}\hrule\vspace{1mm}


\section{第 3 章 逐项数值对比}


\subsection{问题 1（1800 s）}


\begin{center}
\small
\begin{tabular}{lllll}
\toprule
量 & 方案 X & 方案 B & 差 & 差异归因 \\
\midrule
中心温度 & 34.0425 & 33.5753 & $+0.4672$ & 环境 \\
表面温度 & 37.1989 & 36.7856 & $+0.4133$ & 环境 \\
中心含水 & 2.5500 & 2.5500 & $0$ & — \\
表面含水 & 1.5109 & 1.5102 & $+0.0007$ & 环境（$C_{air}$ 拟合值偏高 $0.0015$） \\
\bottomrule
\end{tabular}
\end{center}


\subsection{问题 2（3 h）}


\begin{center}
\small
\begin{tabular}{lllll}
\toprule
量 & 方案 X & 方案 B & 差 & 归因 \\
\midrule
中心温度 & 49.9051 & 49.8495 & $+0.0556$ & 环境（温差随温度场趋于均匀而收敛） \\
中心含水 & 1.7673 & 1.7662 & $+0.0011$ & 环境 \\
表面含水 & 1.0083 & 1.0081 & $+0.0002$ & 环境 \\
\bottomrule
\end{tabular}
\end{center}


\subsection{问题 3、4 的烘干时长}


\begin{center}
\small
\begin{tabular}{llll}
\toprule
量 & 方案 X & 方案 B & 差 \\
\midrule
问题 3 $t_{dry}$ & 205 559 s = \textbf{57.10 h} & 205 809.78 s = \textbf{57.1694 h} & $+0.122\%$ \\
问题 4 $t_{dry}$ & 182 765 s = \textbf{50.77 h} & 182 954.56 s = \textbf{50.8207 h} & $+0.104\%$ \\
\bottomrule
\end{tabular}
\end{center}


\textbf{关键印证}：方案 X 自己的不确定度表里写着


\begin{quote}
「边界条件（用噪声数据）：$+0.06$ h（$+0.11\%$）」
\end{quote}


$57.10+0.06=57.16$ h —— \textbf{与方案 B 的 57.1694 h 吻合到 0.01 h}。


\textbf{所以方案 B 的值恰好就是方案 X 的「改用原始噪声数据」变体。} 两篇论文在这一点上其实是\textbf{互相印证的}，只是各自选择了不同的基准。


\subsection{问题 4 的不收缩对照：方案 X 更完整}


\begin{center}
\small
\begin{tabular}{llll}
\toprule
 & 方案 X & 方案 B & 我的独立复算 \\
\midrule
附录 4 物性 + $R\equiv2$ cm & \textbf{128.90 h} & \textbf{未报告} & 129.09 h（用方案 B 求解器 + 分段线性环境） \\
附录 4 物性 + $R(t)$ & 50.77 h & 50.8207 h & 50.8207 h ✓ \\
\bottomrule
\end{tabular}
\end{center}


\textbf{这是两篇论文最重要的一处方法学差别。}


\textbf{方案 X 的做法}：做对照实验，用\textbf{同一张物性表}只切换半径，从而把两个效应分开：


\begin{equation*}
57.17\ \text{h}\xrightarrow[\text{物性表}]{\times2.254(+125.4\%)}128.90\ \text{h}\xrightarrow[\text{收缩}]{\times0.3939(\mathbf{-60.6\%})}50.77\ \text{h}
\end{equation*}


方案 X 在论文的「模型优点」里明确写着「收缩使干燥缩短 \textbf{60.6\%}」 ✓


\textbf{方案 B 的做法}：直接比较「问题 3（附录 3，无收缩，57.17 h）」与「问题 4（附录 4，有收缩，50.82 h）」，写道


\begin{quote}
「收缩使半径在 40 h 内由 2.0 cm 降到约 1.2 cm，扩散路径缩短约 40\%，同时表面积减小使外边界通量下降，\textbf{净效果是比问题 3 缩短 6.35 h（11\%）}。」
\end{quote}


\textbf{这句话是错的}：它把物性表效应（$+125.4\%$）与收缩效应（$-60.6\%$）两个方向相反、量级都远大于 11\% 的效应混在一起，净剩 $-11\%$，然后把这个净值当成了「收缩的效果」。


\textbf{真实情况}：收缩使烘干时间缩短 \textbf{60.6\%}（从 128.90 h 降到 50.77 h），与 $R^2$ 从 4.00 降到 $1.44\ \mathrm{cm^2}$（即 0.36 倍）的量级完全一致。


\textbf{结论}：方案 X 在问题 4 的物理归因上\textbf{完全正确且更完整}；方案 B 在这一处\textbf{实质性错误}，且因此漏掉了一个本可极大增强说服力的对照实验。


\vspace{1mm}\hrule\vspace{1mm}


\section{第 4 章 三处模型形式歧义（双方都存在）}


\subsection{干基密度的梯度项 $\nabla\ln\rho_d$}


\textbf{理论}：以干基浓度 $C$ 为变量、水分通量为 $\rho_dD\nabla C$ 时，严格的质量守恒给出


\begin{equation*}
\frac{\mathrm DC}{\mathrm Dt}=\nabla\!\cdot\!(D\nabla C)+D\nabla C\!\cdot\!\nabla\ln\rho_d
\end{equation*}


其中 $\rho_d=\rho(C)/(1+C)$。\textbf{第二项在两篇论文的基准模型中都被略去}。


\textbf{两篇论文的处理}：


\begin{center}
\small
\begin{tabular}{lll}
\toprule
 & 是否讨论 & 是否量化 \\
\midrule
方案 X & 未提及 & 未量化 \\
方案 B & §7.2(1) 明确讨论 & **量化：$-3.33\%$** \\
\bottomrule
\end{tabular}
\end{center}


\textbf{方案 B 的量化值（我复现）}：问题 3 的 $t_{dry}$ 从 205 809.78 s（57.1694 h）变为 198 956.60 s（55.2657 h），$−3.33\%$。


\textbf{怎么看待这个歧义}：题面给出的边界条件是 $-D\partial_rC|_R=h_m(C_s-C_{air})$——用的是**不加 $\rho_d$ 的通量**。如果内部加上 $\rho_d$、边界不加，反而前后不一致。所以「略去该项」是\textbf{与题面边界条件自洽}的读法，两篇论文的选择都可以接受；但\textbf{都应把它登记为不确定度}。方案 B 做到了，方案 X 没有。


\begin{quote}
\textbf{一个连带发现}：方案 B 的论文把 $-3.33\%$ 归给了它的式 (13)「严格守恒形式」。但我把式 (13) 完整展开（有效容量 $\rho_d+C\rho_d'$）后重算，得到的是 **191 190.25 s = 53.1084 h（$−7.10\%$）**。也就是说，$-3.33\%$ 对应的其实是「略去 $C\partial_t\rho_d$」的近似模型，不是式 (13)。这导致方案 B 摘要中的模型形式不确定度区间 \textbf{55.3—57.2 h} 偏窄，应为 \textbf{53.1—57.2 h}。
\end{quote}


\subsection{几何相似收缩与干物质守恒的内在张力}


方案 X 的物质坐标推导（原文 5.5 节）非常漂亮，它得到两条结论：


\begin{enumerate}
\item 干物质守恒给出 $\rho_d(\xi,t)R(t)^2=\rho_d(\xi,0)R_0^2$（材料不变量）；
\item 初始含水率均匀 ⟹ $\rho_d(\xi,0)$ 与 $\xi$ 无关 ⟹ **$\rho_d$ 全程与 $\xi$ 无关**，从而可以在方程中约去。
\end{enumerate}


\textbf{但这两条与最终结果之间存在一个逻辑缺口}：


\begin{equation*}
\rho_d=\frac{\rho(C)}{1+C}\ \text{与}\ \xi\ \text{无关}\ \Longrightarrow\ \frac{\rho(C)}{1+C}\ \text{处处相等}
\end{equation*}


而 $\rho(C)/(1+C)$ 是 $C$ 的\textbf{严格单调函数}（附录 4：导数 $-670/(1+C)^2<0$），所以「处处相等」只能推出 **$C$ 处处相等**——这与实际算出的 $C$ 沿半径显著变化（如 3 h 时中心 1.7673、表面 1.0083）\textbf{直接矛盾}。


\textbf{根源}：严格的干物质守恒 + 几何相似收缩（$r=\xi R(t)$）+ $\rho_d=\rho(C)/(1+C)$ 这三条\textbf{不能同时成立}。若坚持 $r=\xi R(t)$，就必须放弃「$\rho_d$ 由局部 $C$ 决定」。


\textbf{正确的写法}（我在详解报告中给出）：不要用「干物质守恒」来论证，而应\textbf{直接把「几何相似收缩」作为一条假设}，然后用「对流项与坐标变换项精确抵消」这个\textbf{恒等式}完成推导——\textbf{最终得到的方程完全一样，但只需要一条假设，没有逻辑缺口}。


方案 B 在这一点上处理得更明确：它独立建立了非均匀收缩自洽模型（M1），正面处理「局部失水多的位置收缩更多」这一可能，并量化了它与均匀仿射假设的差异（$-6.3\%$）。\textbf{这是方案 B 的真实方法学贡献}，尽管其 M1 的边界系数口径本身还有争议（详见我在上一轮审阅报告中的分析）。


\subsection{蒸发潜热}


\begin{center}
\small
\begin{tabular}{lll}
\toprule
 & 方案 X & 方案 B \\
\midrule
基准是否含潜热 & 不含 & 不含 \\
理由 & 题给 $(h,h_m)$ 与热质比拟的 Lewis 关系\textbf{相差 4 个数量级}（$3.13\times10^7$ vs $1.10\times10^3$），真实质量通量相差 250 倍、无法判定 & 能量审计：60 h 内供热仅为潜热需求的 $1/15.7$ \\
潜热变体算例 & 做了一次，\textbf{数值发散}，如实报告「不可用」 & 报告「若显式加入，药材会被冷却到环境温度以下」 \\
\bottomrule
\end{tabular}
\end{center}


\textbf{两篇论文的结论一致}：题面数据不支持一致地计入潜热。方案 X 的分析更进一层——它指出了\textbf{题面参数本身的不自洽}（Lewis 关系），这比「能量不够」的论证更本质。


\textbf{两者的共同优点}：都没有给出一个来自未收敛算例的数字。这是诚实的做法。


\vspace{1mm}\hrule\vspace{1mm}


\section{第 5 章 方法学逐项对比}


\subsection{方案 X 更好的地方}


\begin{center}
\small
\begin{tabular}{ll}
\toprule
项 & 说明 \\
\midrule
\textbf{问题 4 的收缩归因} & 做对照实验，得出 $-60.6\%$（正确）；方案 B 写成 $-11\%$（错误） \\
\textbf{环境口径的透明度} & 明确报告「拟合 vs 实测」两种口径的差异（最大 0.52 ℃、$t_{dry}$ 0.11\%），并计入不确定度表；方案 B 只对 $t_{dry}$ 做了插值方式对比，未报告对表 1 的影响 \\
\textbf{终止判据的严谨性} & 明确区分「全域最大值」与「$\arg\max$ 下标」，并指出下标在剖面均匀时是浮点并列、不携带信息；方案 B 只在末态检查了单调性 \\
\textbf{解析解验证} & 与 Bessel 级数解对比（全场 $1.0\times10^{-3}$ K、中心 $1.2\times10^{-5}$ K），并给出\textbf{观测收敛阶}（空间 2.1、时间 1.9）；方案 B 只有内部一致性检验与制造解 \\
\textbf{外部数据交叉校验} & 发现附件 2 与附录 4 在干基守恒下最大不自洽 17.8\%，并给出「无解」的物理含义；方案 B 也有类似分析但未做定量上界 \\
\textbf{潜热问题的深度} & 指出 $(h,h_m)$ 与 Lewis 关系差 4 个数量级，说明 $h_m$ 是模型内部约定 \\
\textbf{多实现互证} & 节点中心 FV / 单元中心 FV / 直线法三种实现，$0.0066\%$ \\
\textbf{参考文献} & 有真实引用（Crank、Incropera、Luikov、Adrover、Mujumdar、Karatas 等）；方案 B 的 7 条文献\textbf{正文中一次都没引用} \\
\bottomrule
\end{tabular}
\end{center}


\subsection{方案 B 更好的地方}


\begin{center}
\small
\begin{tabular}{ll}
\toprule
项 & 说明 \\
\midrule
\textbf{端面效应的实验设计} & 二维轴对称 + 「同网格只切换 $z=L/2$ 边界」的隔离实验（$-0.004\%$）。方案 X 只有量级论证（轴向/径向特征时间相差 156 倍），虽然结论正确但没有直接验证 \\
\textbf{终止时刻的定位精度} & 连续事件函数 + Brent 求根（精度 $10^{-9}$）优于在离散网格上线性插值。不过方案 X 的分辨率分析显示后者误差仅 0.02\%，实际影响很小 \\
\textbf{非均匀收缩的独立建模} & M1 自洽模型正面处理了 4.2 节的逻辑缺口，并量化 $-6.3\%$；方案 X 未涉及 \\
\textbf{公式判读的风险提示} & 方案 B 把「$e^{-0.89/C}$ 是分式不是乘积」这一判读风险作为全文重点，并保留原版面截图作证据。（两篇论文的实际读法都正确，但方案 B 的\textbf{风险意识与证据留存}更值得借鉴） \\
\textbf{规范符合性} & 方案 B 有 216 格的跨实现逐格对拍表；方案 X 未见同类逐格对拍 \\
\bottomrule
\end{tabular}
\end{center}


\subsection{两篇论文各自的问题}


\begin{center}
\small
\begin{tabular}{lll}
\toprule
 & 方案 X & 方案 B \\
\midrule
严重 & ① 几何相似收缩推导中的逻辑缺口（4.2 节）；② 表 1 采用拟合边界使其系统性偏高至 0.52 ℃（已披露） & ① 式 (13) 归属错误，导致不确定度区间偏窄；② 式 (10) 温度守恒式在问题 2/3/4 下不成立；③ W1 单元测试第一张表空转 \\
中等 & ③ 未登记干基梯度项的不确定度；④ 变体表「问题3 基准 57.08 h」与其他处的 57.10 h 不一致 & ④ 问题 4 收缩效应归因错误（$-11\%$ vs $-60.6\%$）；⑤ 参考文献 0 次引用；⑥ $\delta$ 数值标签错、弹性差 100 倍、表 D 三处自相矛盾 \\
轻微 & ⑤ 潜热变体数值发散（已如实报告，属实现问题而非结论问题） & ⑦ 若干定性叙述的数量级失准（水分扩散长度、前沿位置等） \\
\bottomrule
\end{tabular}
\end{center}


\vspace{1mm}\hrule\vspace{1mm}


\section{第 6 章 结论与建议}


\subsection{对两套结果的最终判定}


\begin{center}
\small
\begin{tabular}{ll}
\toprule
问题 & 判定 \\
\midrule
问题 1、2 的温度 & \textbf{两者都对，差异全部来自环境口径}。若题面要求「按附件 1 的数据」，则分段线性（方案 B 的值）更忠实；若允许对噪声数据做物理光滑，则方案 X 的值成立。\textbf{两篇都应把对方的值作为对照报告出来。} \\
问题 1、2 的含水率 & 两者一致（差 $\le0.06\%$），且已被两种独立求解器互证 \\
问题 3 的烘干时长 & \textbf{57.10 h（X）与 57.17 h（B）相差 0.12\%，完全由环境口径解释}。建议表述为「\textbf{57.1$\pm$0.1 h（环境口径）}」 \\
问题 4 的烘干时长 & 50.77 h（X）与 50.82 h（B），差 0.10\%，同样由环境口径解释 \\
\textbf{问题 4 的收缩效应} & **方案 X 正确（$-60.6\%$），方案 B 错误（$-11\%$）。** 这是本轮比对最重要的结论 \\
\bottomrule
\end{tabular}
\end{center}


\subsection{给方案 X 的改进建议}


\begin{enumerate}
\item \textbf{修正 5.5 节的推导}：删去「由干物质守恒推出 $\rho_d$ 与 $\xi$ 无关」这一步（它与非均匀剖面矛盾），改为「\textbf{假设几何相似收缩} + 对流项精确抵消」的推法。最终方程不变，但逻辑严密。
\item \textbf{补充干基梯度项的量化}：把 $\nabla\ln\rho_d$ 项加回去，给出 $t_{dry}$ 的变化（参考值 $-3.33\%$），并说明为什么基准模型选择略去它（与题面边界条件自洽）。
\item \textbf{在摘要中同时给出两种环境口径的结果}，例如「按实测插值 57.17 h / 按物理拟合 57.10 h」，避免读者只看到其中一个。
\item \textbf{对齐变体表的基准值}（57.08 与 57.10 不一致）。
\item \textbf{补一个二维端面的隔离实验}（成本很低），把「端面可忽略」从量级论证升级为直接验证。
\end{enumerate}


\subsection{给方案 B 的改进建议}


\begin{enumerate}
\item \textbf{修正问题 4 的收缩归因}：补不收缩对照（128.90 h），把「收缩使时长缩短 11\%」改为「**物性表效应 $+125.4\%$、收缩效应 $-60.6\%$、净 $-11\%$**」。
\item \textbf{修正式 (13) 的归属}：把「严格守恒形式」的实现改为有效容量 $\rho_d+C\rho_d'$，或把表述改为「略去 $C\partial_t\rho_d$ 的近似模型」；摘要的模型形式不确定度区间由 \textbf{55.3—57.2 h} 扩为 \textbf{53.1—57.2 h}。
\item \textbf{修正式 (10) 的温度守恒式}：问题 2/3/4 下应写作 $\int\sigma\rho c_p\partial_tT\,\mathrm d\sigma=-\frac hR(T_s-T_\infty)$。
\item \textbf{修复 W1 单元测试并落盘日志}。
\item \textbf{补全参考文献引用}（当前正文中 cite 命令出现 0 次）。
\item \textbf{补报告「拟合 vs 插值」两种环境口径对表 1 的影响}（方案 X 已经做了，这 0.52 ℃ 的偏差是应该披露的）。
\end{enumerate}


\subsection{一个可供两篇共用的表述模板}


\begin{quote}
本题的烘房环境由附件 1 的\textbf{带噪实测序列}给出，把它转换为连续边界函数有两种自洽做法： （i）分段线性插值——精确通过每个实测点，不做额外假设； （ii）一阶惯性式拟合——物理上与烘房升温过程一致，但残差滞后 1 阶自相关达 0.83，属模型形式失配，在 $60\sim2300$ s 内系统性高于实测达 0.96 ℃。 两者使问题 1 的表 1 相差最多 0.52 ℃，使问题 3 的烘干时长相差 0.11\%。 本文正文采用（\_\_），并给出（\_\_）的结果作为对照。
\end{quote}


\textbf{这一段本身就是一篇优秀论文该有的样子}：不只给答案，还给「答案对什么敏感」。


\subsection{本轮比对对方案 B 审阅结论的影响}


\begin{center}
\small
\begin{tabular}{ll}
\toprule
项 & 变化 \\
\midrule
方案 B 的结果文件（表 1—表 6、\texttt{result1}–\texttt{result4}） & \textbf{无需修改}（自身内部一致，且已被另一套独立求解器在四位小数上复现） \\
新增 1 条应采纳的对照实验 & 问题 4 不收缩对照 128.90 h \\
新增 1 处必须更正的归因 & 收缩效应 $-11\%$ $\rightarrow$ **$-60.6\%$** \\
新增 1 条应登记的不确定度 & 环境口径，对 $t_{dry}$ 影响 $0.11\%$ \\
原审阅报告的 A 级问题 & 不变（4 项：式 13 归属、式 10 温度式、W1 空转、V1/V2 证据链） \\
\bottomrule
\end{tabular}
\end{center}

